\documentclass[10pt,a4paper]{amsart}
\usepackage[margin=19mm]{geometry}
\usepackage{amsmath,amssymb,mathtools}
\usepackage[scr=rsfs,cal=euler]{mathalfa}
\usepackage{xfrac}
\usepackage[unicode,hidelinks,bookmarks=false]{hyperref}
\newcommand{\ie}{i.e.\ }
\newcommand{\cf}{cf.\ }
\newcommand{\resp}{respectively\ }
\newcommand{\Subsection}{\subsection}
\providecommand{\nobreakdash}{\nobreak\hbox{-}}
\setlength{\emergencystretch}{2em}
\multlinegap0pt
\usepackage{amssymb}
\usepackage{mathtools}
\usepackage[shortlabels]{enumitem}
\newtheorem{lemma}{Lemma}
\DeclareMathOperator{\Ad}{Ad}
\newcommand{\HS}{{\rm HS}}
\newcommand{\La}{\lambda}
\newcommand{\ol}[1]{\overline{#1}}
\title{$L^2$-cohomology and deformations of \\ the left regular representation}
\author{Christoph Gamm, Andreas Thom}
\date{Source: arXiv:2607.17670v1 (CC BY 4.0)}
\begin{document}
\maketitle

\noindent\emph{Source and licence.} $L^2$-cohomology and deformations of \\ the left regular representation. Version arXiv:2607.17670v1 (CC BY 4.0), distributed by arXiv under the Creative Commons Attribution 4.0 International licence, http://creativecommons.org/licenses/by/4.0/, as recorded in the arXiv licence metadata for this exact version and shown on the abstract page https://arxiv.org/abs/2607.17670v1. Full licence notice: \texttt{licenses/arxiv-2607.17670-CC-BY-4.0.txt}.\par\smallskip

\noindent\textit{Editorial scope.} Counted unit: the article's lemma lem:transfer (source0.tex lines 370--377). The article's complete original proof of it is delivered with it (source0.tex lines 379--493, one \texttt{proof} environment of 115 lines). No mathematics is written, repaired or re-derived: the statement and the proof are the article's own lines, in the article's own notation, and no retained line is substituted or re-flowed.  Retained uncounted as the source-local context the proof needs: lines 296--307; lines 348--354.  Nothing was omitted from any retained span, so the package carries no omission record.  No retained line contains a citation, so the wrapper carries no bibliography and no bibliography entry.  No retained line contains a figure environment, an image-inclusion command or drawing code, so no image is delivered and the package declares no asset dependency.  Editorial formatting only, in addition: an AMS \texttt{amsart} preamble that restates the article's own declarations and macros used by the retained lines, namely the environments \texttt{lemma} on separate counters, together with the standard packages those lines need.  The retained environments print in this document as Lemma 1 in order of appearance, whereas the article prints the same items with the numbering of its own full document.  The complete native source of the article as distributed in the arXiv e-print for this exact version is bundled unchanged as \texttt{sources/205590/source0.tex}.
\begin{lemma}\label{lem:openmapping}
Let $V$ be a Fr\'echet $G$-module and let $n\ge 0$. Assume that $H^{n+1}(G,V)=0$. Then the induced map
\[
\bar d_n\colon C^n(G,V)/Z^n(G,V)\to Z^{n+1}(G,V)
\]
is an isomorphism of Fr\'echet spaces. In particular, for every finite set $E\subset G^n$ there exist a finite set $F\subset G^{n+1}$ and a constant $\kappa>0$ such that every cocycle $z\in Z^{n+1}(G,V)$ admits a primitive $c\in C^n(G,V)$ with
\[
d_nc=z,
\qquad
p_E(c)\le \kappa p_F(z).
\]
\end{lemma}

\begin{lemma}[Fell]\label{lem:HSdecomp}
There is an isometric isomorphism of Hilbert $G$-modules
\[
\HS(\ell^2 G)_{\Ad\La}\cong \ell^2(G)\otimes \ol{\ell^2(G)}\cong \ell^2(G)\otimes \mathcal K,
\]
where $\mathcal K$ is a separable Hilbert space with trivial $G$-action.
\end{lemma}

\begin{lemma}\label{lem:transfer}
Let $G$ be finitely generated. Then
\[
H^2\bigl(G,\HS(\ell^2 G)_{\Ad\La}\bigr)=0
\quad\Longleftrightarrow\quad
H^2\bigl(G,\ell^2(G)\bigr)=0.
\]
\end{lemma}

\begin{proof}
By Lemma~\ref{lem:HSdecomp}, it suffices to compare
$\ell^2(G)\otimes\mathcal K$ with $\ell^2(G)$, where
$\mathcal K$ has trivial action. Fix an orthonormal basis
$(e_j)_{j\ge 1}$ of $\mathcal K$.

Assume first that
$
H^2\bigl(G,\ell^2(G)\bigr)=0.
$
Let
$
c\in Z^2\bigl(G,\ell^2(G)\otimes\mathcal K\bigr)
$
and write
\[
c(g,h)=\sum_{j\ge 1}c_j(g,h)\otimes e_j.
\]
Then each $c_j$ belongs to
$Z^2\bigl(G,\ell^2(G)\bigr)$.

Fix a finite symmetric generating set $S$ of $G$. By
Lemma~\ref{lem:openmapping}, there exist a finite set
$F\subseteq G^2$ and a constant $\kappa>0$ such that, for
every $j$, there is a primitive
$
b_j\in C^1\bigl(G,\ell^2(G)\bigr),$ with $
d_1b_j=c_j,$
satisfying
\[
\max_{s\in S}\|b_j(s)\|_2
\leq
\kappa\max_{(g,h)\in F}\|c_j(g,h)\|_2.
\]
Consequently, for every $s\in S$,
\[
\begin{aligned}
\sum_{j\ge 1}\|b_j(s)\|_2^2
&\leq
\kappa^2
\sum_{j\ge 1}
\max_{(g,h)\in F}\|c_j(g,h)\|_2^2 \\
&\leq
\kappa^2
\sum_{(g,h)\in F}
\sum_{j\ge 1}\|c_j(g,h)\|_2^2 
=
\kappa^2
\sum_{(g,h)\in F}\|c(g,h)\|_2^2
<\infty.
\end{aligned}
\]
Thus
$
b(s):=\sum_{j\ge 1}b_j(s)\otimes e_j
$
is well defined for every $s\in S$.

Now let $g\in G$ and choose a word
$
g=s_1\cdots s_r,$ with $
s_1,\ldots,s_r\in S.$
Since $d_1b_j=c_j$, one has
\[
b_j(g)
=
b_j(s_1)
+
\sum_{i=2}^r
s_1\cdots s_{i-1}\cdot b_j(s_i)
-
\sum_{i=2}^r
c_j(s_1\cdots s_{i-1},s_i).
\]
It follows that
\[
\begin{aligned}
b(g)
:={}&
b(s_1)
+
\sum_{i=2}^r
s_1\cdots s_{i-1}\cdot b(s_i)
-
\sum_{i=2}^r
c(s_1\cdots s_{i-1},s_i)
\end{aligned}
\]
is a well-defined element of
$\ell^2(G)\otimes\mathcal K$. Coordinatewise, it equals
$
\sum_{j\ge 1}b_j(g)\otimes e_j.
$
Hence
$
b\in C^1\bigl(G,\ell^2(G)\otimes\mathcal K\bigr),
$
and the identity $d_1b=c$ follows coordinatewise. Therefore
$
H^2\bigl(G,\ell^2(G)\otimes\mathcal K\bigr)=0.
$

Conversely, choose a basis vector $e_1\in\mathcal K$. The
equivariant maps
$
\iota(v)=v\otimes e_1
$
and
$
P\left(\sum_jv_j\otimes e_j\right)=v_1
$
satisfy $P\circ\iota={\rm id}$. Passing to cohomology shows that
vanishing for $\ell^2(G)\otimes\mathcal K$ implies vanishing
for $\ell^2(G)$.
\end{proof}

\end{document}
